\section{模型假设}

为简化问题并建立可解的数学模型，提出以下假设：

{\raggedright
\begin{enumerate}
\item \textbf{样本独立同分布。} 结直肠癌（Zeller）、炎症性肠病（metahit）和肥胖症（Chatelier）三个数据集内部的样本，各自视为独立同分布抽样；三个数据集分别建模，不混合使用。因为不同研究的数据存在批次差异，直接混合可能引入伪信号。
\item \textbf{相对丰度为定和成分数据。} 每个样本的物种级相对丰度之和约为 100，存在定和约束。因此在线性建模前必须做 CLR 变换，解除“此消彼长”带来的伪相关；前期验证显示，CLR 变换可使模型性能提升约 0.10--0.15。
\item \textbf{零值视为“未检出”的真实稀疏。} 原始矩阵中 92.21\% 的值为零，这些零值视为微生物未被检出，而不是真实缺失。对零值采用乘法替换伪值 $\delta=6.5\times10^{-6}$，该值约为检出限的 0.65 倍，且结果对 $\delta$ 的微小扰动不敏感。
\item \textbf{近全零特征不含判别信息。} 零值占比超过 95\% 的特征几乎不携带判别信息，因此从 1331 维中剔除 1067 维，保留 264 维。这些特征在后续筛选中的选中率接近 0，剔除不会损失判别能力。
\item \textbf{标签可靠。} 患病/健康标签由数据来源给定，本文认为标签是可靠的。Zeller 数据中的 small adenoma 样本按主方案归入健康组，该组样本量最大；其余分组方式作为敏感性分析。
\item \textbf{批次效应不主导全局结构。} 三个数据集虽然来自不同研究/平台，但初步分析显示样本的 silhouette 系数为 0.070，接近 0，说明样本并没有按研究来源明显分群。跨疾病性能差异主要来自疾病特异信号与患病定义差异。
\item \textbf{判别信号以“存在/缺失”为主、“丰度高低”为辅。} 数据分析显示，物种“有/无”差异与 AUC 的相关性（0.86--0.89）明显高于非零丰度差异（0.63--0.67），因此本文分别用“存在/缺失”和“非零丰度”两路检验来捕捉不同信号是合理的。
\item \textbf{无外部验证数据。} 模型泛化能力通过内部交叉验证（分层 5 折交叉验证 + LOOCV）和跨疾病评估方案（LODO）检验；全量训练集上的 AUC 只能作为乐观上界，不能用来宣称模型性能。
\end{enumerate}
}
\fussy

\noindent\textbf{敏感性提示}：模型对“样本独立同分布”和“无批次效应”假设的违背比较敏感。如果三个数据集之间存在较强的批次效应，跨疾病预测性能可能会进一步下降。正文敏感性分析已经对 small adenoma 的分组方式、$\tau$ 阈值和第三问阈值偏移做了量化检验。
